\documentclass[10pt, a4paper]{article}

% Parametri che modificano il file main.tex
% Le uniche parti da cambiare su main.tex sono:
% - tabella versioni
% - testo file

\def\titolo{Analisi dei \\ \ requisiti} % \\ per andare a capo, \ per spazio
\def\titoloHeader{Analisi dei requisiti}
\def\data{v2.0.0(0)}

% \def\listaComponenti{
% Bresolin G.,
% Campese M.,
% Ciriolo I.,
% Dugo A.,
% Feltrin E.,
% Michelon R.,
% Orlandi G.
% }

\usepackage{style}
\usepackage{headerfooter}
\usepackage{xhfill}



\title{\titolo}
\author{SWEetCode}
\renewcommand{\thesection}{}
\renewcommand{\thesubsection}{}
%\renewcommand{\thesubsubsection}{}
\newcommand{\pg}{\textsubscript{|g|}}
\usepackage{titlesec}
\titleformat{\section}
  {\normalfont\fontsize{28}{15}\bfseries}{\thesection}{1em}{}


\begin{document}

% PRIMA PAGINA
\begin{titlepage}
    \thispagestyle{empty}
    \begin{tikzpicture}[remember picture, overlay]
        % TRIANGOLI
        \draw[fill=secondarycolor, secondarycolor] (current page.north west) -- (current page.south west) -- (8.8, -28);
        \draw[fill=primarycolor, primarycolor] (-3, 5) -- (4, -13.6) -- (11, 5);

        % LOGO
        \node [xshift=-5cm, yshift=25cm] (logo) at (current page.south east) {\includegraphics[width=6.5cm]{img/logo.png}};

        % SWEETCODE - DATE
        \node [anchor=north east, align=right, xshift=-1.2cm, yshift=20.5cm, text=black] (sweetcode) at (current page.south east) {\fontsize{32pt}{36pt}\selectfont SWEetCode};
        \draw[line width=4pt, lightcol] ([xshift=-3cm, yshift=-0.37cm]sweetcode.south west) -- ([yshift=-0.37cm]sweetcode.south east);
        \node [anchor=north east, align=right, xshift=-1.2cm, yshift=18.7cm, text=black] (date) at (current page.south east){\fontsize{24pt}{24pt} \selectfont Verbale \tipoVerb};

        % NOME FILE
        \node [anchor=north east, text width=15cm, align=right, xshift=-1.2cm, yshift=17cm, text=black] (titolo) at (current page.south east){\fontsize{48pt}{48pt}\textbf{\data}};

        % BOX DATI PARTECIPANTI
        \ifthenelse{\equal{\tipoVerb}{Esterno}}{
            \node[anchor=north east, xshift=-1.2cm, yshift=14.5cm, minimum width=8cm] (box) at (current page.south east){};
        }{
            \node[anchor=north east, xshift=-1.2cm, yshift=12.5cm, minimum width=8cm] (box) at (current page.south east){};
        }

        % RESPONSABILE
        \node[anchor=north west, align=left] (dati2) at (box.north west) {\fontsize{15pt}{15pt}\selectfont \textbf{Responsabile}};
        \draw[line width=4pt, lightcol] (dati2.south west) -- ([xshift=8cm]dati2.south west);
        \node[anchor=north west, align=left] (dati21) at (dati2.south west){\fontsize{13pt}{13pt}\selectfont \nomeResp};

        % VERIFICATORE
        \node[anchor=north west, yshift=-1cm, align=left] (dati3) at (dati21.north west) {\fontsize{15pt}{15pt}\selectfont \textbf{Verificatore}};
        \draw[line width=4pt, lightcol] (dati3.south west) -- ([xshift=8cm]dati3.south west);
        \node[anchor=north west, align=left] (dati31) at (dati3.south west){\fontsize{13pt}{13pt}\selectfont \nomeVer};

        % SEGRETARIO DI RIUNIONE
        \node[anchor=north west, yshift=-1cm, align=left] (dati4) at (dati31.north west) {\fontsize{15pt}{15pt}\selectfont \textbf{Segretario di Riunione}};
        \draw[line width=4pt, lightcol] (dati4.south west) -- ([xshift=8cm]dati4.south west);
        \node[anchor=north west, align=left] (dati41) at (dati4.south west){\fontsize{13pt}{13pt}\selectfont \nomeSegr};
        
        % UNIPD - SWE
        \node [xshift=4.4cm, yshift=2.3cm, draw, secondarycolor, text=white] (uni) at (current page.south west) {\fontsize{20pt}{20pt} \selectfont Università di Padova};
        \node [xshift=0.65cm, yshift=0.7cm, draw, secondarycolor, text=white, below=of uni] (corso) {\fontsize{20pt}{20pt}\selectfont Ingegneria del Software};

        % FIRMA AZIENDA
        \ifthenelse{\equal{\tipoVerb}{Esterno}}{
            \draw[line width=4pt, lightcol] ([xshift=-1.2cm, yshift=7cm]current page.south east) -- ([xshift=-8cm, yshift=7cm]current page.south east);
            \node [xshift=-4.8cm, yshift=8cm] (logo) at (current page.south east) {\includegraphics[width=6cm]{img/firme/\firmaAzienda}};
            \node[anchor=north west, xshift=12.9cm, yshift=6.65cm, align=left] at (current page.south west)
            {\fontsize{13pt}{13pt}\selectfont L'azienda: \nomeAzienda};
        }

        % FIRMA CODING COWBOYS
        \draw[line width=4pt, lightcol] ([xshift=-1.2cm, yshift=4.35cm]current page.south east) -- ([xshift=-8cm, yshift=4.35cm]current page.south east);
        \node [xshift=-4.8cm, yshift=5.5cm] (logo) at (current page.south east) {\includegraphics[width=6cm]{img/firme/menon.png}};
        \node[anchor=north west, xshift=12.9cm, yshift=4cm, align=left] at (current page.south west)
        {\fontsize{13pt}{13pt}\selectfont Coding Cowboys: Menon G.};

        % FIRMA
        \draw[line width=4pt, lightcol] ([xshift=-1.2cm, yshift=1.8cm]current page.south east) -- ([xshift=-8cm, yshift=1.8cm]current page.south east);
        \node [xshift=-4.8cm, yshift=2.45cm] (logo) at (current page.south east) {\includegraphics[width=6cm]{img/firme/\firmaResp}};
        \node[anchor=north west, xshift=12.9cm, yshift=1.45cm, align=left] at (current page.south west)
        {\fontsize{13pt}{13pt}\selectfont Il Responsabile: \nomeResp};
        
    \end{tikzpicture}
\end{titlepage}

% REGISTRO DELLE VERSIONI
%Registro in ordine dalla più recente alla meno recente!
{\renewcommand{\arraystretch}{1.5}
\section*{Registro delle versioni}

\begin{xltabular}{\textwidth}{c|c|c|c|X}
\label{tab:long}

\textbf{Versione} & \textbf{Data} & \quantities{\textbf{Responsabile di}\\\textbf{stesura}}& \textbf{Revisore} & \quantities{\textbf{Dettaglio e}\\\textbf{motivazioni}} \\
\endfirsthead

\textbf{Versione} & \textbf{Data} & \quantities{\textbf{Responsabile di}\\\textbf{stesura}}& \textbf{Revisore} & \quantities{\textbf{Dettaglio e}\\\textbf{motivazioni}} \\
\endhead

\multicolumn{5}{r}{{Continua nella pagina successiva}} \\
\endfoot

\endlastfoot

\hline
v1.10.0(7) & $2023-12-29$ & \quantities{Michelon R.} & Dugo A. & Aggiunti diagrammi use case e di attività.\\
\hline
v1.10.0(6) & $2023-12-29$ & \quantities{Michelon R.} & Dugo A. & Aggiornamento use case e requisiti.\\
\hline
v1.0.0(14) & $2023-12-06$ & \quantities{Bresolin G.} & Orlandi G. & Aggiornamento UC e requisiti presenti. Introduzione diagrammi dei casi d'uso. Aggiunta di nuovi UC e nuovi requisiti.\\
\hline
v1.0.0(5) & $2023-11-23$ & \quantities{Ciriolo I.} & Campese M. & Introduzione, Descrizione, Studio dei primi Casi d'uso, Studio dei primi requisiti.\\

    
\end{xltabular}}
\newpage

% INDICE
\tableofcontents
\newpage

% INIZIO PAGINE
\section{Introduzione}
\subsection{Scopo del documento}
\paragraph{}Questo documento è stato creato con l'intento di raccogliere tutti i termini che possano risultare sconosciuti ad un lettore esterno o che possano creare ambiguità. Lo scopo è quello di creare una base omogenea di conoscenza specifica. 

\subsection{Impostazione}
\paragraph{}Il documento si sviluppa in una introduzione seguita dalle sezioni organizzate sulle lettere dell'alfabeto; all'interno di queste si troveranno i termini ambigui o specifici del progetto che sono stati contrassegnati nel resto della documentazione con un pedice \textit{|g|} alla fine della parola.



%\subsection{Riferimenti informativi}


%\newpage
\section{A}
\vspace{2em}
\paragraph{Agile}\noindent\hrulefill
\paragraph{}Nel campo dello sviluppo software, il metodo Agile fa riferimento ad un modello di sviluppo che promuove la consegna del software al cliente/destinatario in tempi brevi e frequenti in modo da garantire sempre elevata qualità e un perfetto funzionamento.
Altri aspetti caratteristici di questo metodo sono: formazioni di team di sviluppo piccoli, auto-organizzati caratterizzati da una comunicazione informale, sviluppo incrementale e iterativo e coinvolgimento continuo del cliente.

%\paragraph{API}?
\vspace{2em}
\paragraph{Attore}\noindent\hrulefill
\paragraph{}Elemento che rappresenta un soggetto esterno che interagisce con il sistema modellato. Gli attori possono essere utenti, altri sistemi o qualsiasi altra entità che abbia un ruolo nelle interazioni con il sistema modellato.


%\section{B}



\newpage
\section{C}

\vspace{2em}
\paragraph{Capitolato}\noindent\hrulefill\noindent\hrulefill
\paragraph{} È un documento tecnico generalmente collegato ad un contratto d'appalto e definisce delle specifiche tecniche delle opere che andranno eseguite per effetto del contratto stesso.

\vspace{2em}
\paragraph{Ciclo di vita del software}\noindent\hrulefill\noindent\hrulefill
\paragraph{}Insieme degli stati di maturità che un prodotto software assume tra il suo concepimento e il ritiro in conseguenza alle
attività svolte su di esso.

\vspace{2em}
\paragraph{Committente}\noindent\hrulefill
\paragraph{}Committente è una parola italiana derivante dalla parola latina \textit{committo}. Essa indica una persona o ente che commissiona o ordina un lavoro o un servizio a un fornitore o a un prestatore di
servizi. Nel corso del progetto in questo caso ci si riferisce al Professor Tullio Vardanega.

\newpage
\section{D}

\vspace{2em}
\paragraph{Diagramma di Gantt}\noindent\hrulefill
\paragraph{}Il diagramma di Gantt è costruito partendo da un asse orizzontale a rappresentazione dell’arco temporale
totale del progetto, suddiviso in fasi incrementali (e.g. giorni, settimane, mesi) e da un asse
verticale a rappresentazione delle mansioni o attività che compongono il progetto.


\newpage
\section{E}
\vspace{2em}
\paragraph{Efficacia}\noindent\hrulefill
\paragraph{}L'efficacia indica la capacità di raggiungere l'obiettivo prefissato.

\vspace{2em}
\paragraph{Efficienza}\noindent\hrulefill
\paragraph{}La capacità di azione o di produzione con la massima efficacia ma con il minimo scarto/spesa di risorse e di tempo impiegati.

%\vspace{2em}
%\section{F}

\newpage
\section{G}
\vspace{2em}
\paragraph{Git}\noindent\hrulefill
\paragraph{}Sistema software di controllo di versione distribuito.

\vspace{2em}
\paragraph{GitHub}\noindent\hrulefill
\paragraph{}È un servizio web di hosting specializzato per lo sviluppo di progetti software 
che usa il sistema di controllo di versione \textit{Git\pg}.
%GitHub offre diversi piani per repository privati a pagamento ma anche piani gratuiti molto utilizzati per
%lo sviluppo di progetti open-source.

\vspace{2em}
\paragraph{GitHub Actions}\noindent\hrulefill
\paragraph{}È un servizio di \textit{GitHub\pg} che permette di automatizzare le fasi di build, test e distribuzione nella pipeline di sviluppo in maniera efficiente e scalabile attraverso semplici file di configurazione.

\vspace{2em}
\paragraph{Gulpease}\noindent\hrulefill 
\paragraph{}O “Indice Gulpease” è un indice di leggibilità di un testo tarato sulla
lingua italiana. A differenza di altri ha il vantaggio di utilizzare la lunghezza delle
parole in lettere anziché in sillabe, rendendo più semplice il calcolo automatico.

%\section{H}
\newpage
\section{I}
\vspace{2em}
\paragraph{Ingegneria del Software}\noindent\hrulefill
\paragraph{}L' ingegneria del software è quella disciplina informatica che si occupa dei processi produttivi e delle metodologie di sviluppo finalizzate alla realizzazione di sistemi software.

\vspace{2em}
\paragraph{Issue Tracking System}\noindent\hrulefill
\paragraph{}È un sistema di monitoraggio dei task; è uno strumento che consente di gestire vari casi, di definire la loro priorità e di tenere traccia del relativo avanzamento durante il processo, dall'inizio alla fine. 

\vspace{2em}
\paragraph{ITS}\noindent\hrulefill
\paragraph{}Vedi \textit{Issue Tracking System\pg}

\newpage
\section{J}

\vspace{2em}
\paragraph{Jira}\noindent\hrulefill
\paragraph{}\href{https://www.atlassian.com/it/software/jira}{https://www.atlassian.com/it/software/jira}\\\\
Software utilizzato per l'organizzazione del lavoro tra i membri del team per implementare un \textit{ITS\pg}.




%\section{K}


%\section{L}
\newpage
\section{M}

\vspace{2em}
\paragraph{Milestone}\noindent\hrulefill
\paragraph{}Milestone è il termine letterale per pietra miliare. Figurativamente, la milestone determina importanti traguardi in termini di tempo o di attività; quindi stabiliscono una metrica di avanzamento del progetto.

\vspace{2em}
\paragraph{Minimum Viable Product}\noindent\hrulefill
\paragraph{}L’MVP è una prima versione del prodotto o servizio finale che l’azienda sviluppa. È una versione minima ma funzionante del prodotto o servizio.

\vspace{2em}
\paragraph{MVP}\noindent\hrulefill 
\paragraph{}Vedi \textit{Minimum Viable Product\pg}.

%\section{N}


%\section{O}

\newpage
\section{P}
\vspace{2em}
\paragraph{PDCA}\noindent\hrulefill
\paragraph{}PDCA (acronimo per \textit{Plan-Do-Check-Act}) o \textit{Ciclo di Deming}, è un modello per il
miglioramento continuo della qualità, la quale è rivolta ai processi, non al prodotto.
\begin{itemize}
    \item Plan: definire e pianificare gli obiettivi e i processi necessari per arrivare ai risultati attesi;
    \item Do: realizzare questi obiettivi;
    \item Check: controllare i risultati ottenuti precedentemente;
    \item Act: agire sui processi che non sono stati valutati bene al passo precedente;
\end{itemize}



\vspace{2em}
\paragraph{PoC}\noindent\hrulefill
\paragraph{} Vedi \textit{Proof of Concept\pg}

\vspace{2em}
\paragraph{Processo}\noindent\hrulefill
\paragraph{}Insieme di attività coese che hanno come scopo raggiungere un obiettivo/risultato concreto consumando risorse.

\vspace{2em}
\paragraph{Product Baseline}\noindent\hrulefill
\paragraph{}Seconda revisione di avanzamento prevista dall’attività progettuale. In preparazione a questa revisione il gruppo si dedica alla progettazione architetturale e di dettaglio del prodotto previsto e compreso ad alto livello nella \textit{Requirements and Technology Baseline\pg}.

\vspace{2em}
\paragraph{Proof of Concept}\noindent\hrulefill
\paragraph{}È il prototipo o un esempio di un’idea o di un concetto; questo prototipo dimostra che una particolare soluzione è fattibile attraverso determinate tecnologie e che può essere utilizzata in un ambiente reale. Esso è utile per esplorare le criticità inerenti alle tecnologie e alla loro combinazione.

\vspace{2em}
\paragraph{Proponente}\noindent\hrulefill
\paragraph{}Il proponente è la persona che ha proposto al \textit{Committente\pg} un \textit{Capitolato\pg}
d’appalto.



%\section{Q}
\newpage
\section{R}

\vspace{2em}
\paragraph{Requisito}\noindent\hrulefill
\paragraph{}Una esigenza, funzione o caratteristica di un sistema o delle sue componenti necessarie per il raggiungimento di un obiettivo. 


\vspace{2em}
\paragraph{Requirements and Technology Baseline}\noindent\hrulefill
\paragraph{}Prima revisione di avanzamento di progetto nel corso di \textit{Ingegneria del Software\pg}. Questa revisione richiede di aver realizzato un \textit{PoC\pg} accompagnato dalla documentazione di Analisi dei requisiti.


\newpage
\section{S}
\vspace{2em}
\paragraph{Script}\noindent\hrulefill
\paragraph{}Letteralmente "copione" in inglese, in ambiente informatico rappresenta un programma di dimensioni limitate, solitamente sequenziale e scritto in un linguaggio interpretato. Spesso ha complessità bassa e realizza un singolo task.

\vspace{2em}
\paragraph{Scrum}\noindent\hrulefill
\paragraph{}È un framework  \textit{Agile\pg} per la gestione del ciclo di sviluppo del software che enfatizza il lavoro di squadra, la trasparenza e la responsabilità. Scrum utilizza un metodo incrementale e iterativo per ottimizzare la prevedibilità e il controllo dei rischi.
I ruoli di Scrum nello sviluppo del software sono i seguenti: Product Owner, Scrum Master, Development Team.

\vspace{2em}
\paragraph{Sprint}\noindent\hrulefill
\paragraph{}Periodo di tempo prefissato (14 giorni nel caso del team), entro il quale il gruppo di lavoro si assegna obiettivi, suddivide il lavoro, lo porta a termine e ne verifica la bontà. Questi periodi brevi sono al centro delle metodologie \textit{Agile\pg}.



%\section{T}
\newpage
\section{U}

\vspace{2em}
\paragraph{UML}\noindent\hrulefill
\paragraph{}Acronimo di Unified Modeling Language; È uno standard industriale basato sul paradigma orientato agli oggetti, permette di agevolare la modellazione e progettazione di sistemi informatici.


\vspace{2em}
\paragraph{Use Case}\noindent\hrulefill
\paragraph{}Insiemi di scenari e azioni che hanno in comune uno scopo finale per un \textit{Attore\pg}.

\newpage
\section{V}
\vspace{2em}
\paragraph{Validazione}\noindent\hrulefill
\paragraph{}La validazione è un processo contrattuale ovvero mette in relazione il capitolato con il prodotto. Questo processo ci assicura che quanto prodotto soddisfi le aspettative del cliente finale secondo la corretta individuazione dei requisiti utente.

\vspace{2em}
\paragraph{Verifica}\noindent\hrulefill
\paragraph{}Insieme di attività che si assicura che un prodotto, sistema o servizio soddisfi i requisiti specificati e ne garantisca la qualità, non introducendo errori.


\vspace{2em}
\paragraph{Versionamento}\noindent\hrulefill
\paragraph{}Il versionamento è la gestione di versioni multiple di un insieme di informazioni; il controllo di versione è necessario
per i progetti gestiti da un team collaborativo ed è utilizzato per lo più per gestire la continua evoluzione dei documenti o altro.



\newpage
\section{W}

\vspace{2em}
\paragraph{Way of Working}\noindent\hrulefill
\paragraph{}È un termine per indicare, nel suo complesso, il modo di lavorare adottato (strumenti, norme, tecnologie). Il tutto deve avvenire con procedure determinate, sistematiche e quantificabili.   

%\section{X}
%\section{Y}
%\section{Z}


\end{document}